\documentclass[11pt,a4paper]{article}
\usepackage[margin=24mm]{geometry}
\usepackage{amsmath,amssymb,graphicx,booktabs}
\usepackage{url}
\title{VNAV Lab 3: Modelling and Control of UAVs\\Deliverable 2}
\author{Team Fullcans}
\date{}
\begin{document}
\maketitle

\section*{A. Structure of quadrotors}
We use a right-handed frame with $x$ pointing right and $y$ pointing up in
the top-view exercise figure, and $z$ pointing out of the page. This fixes
all moment signs explicitly. The exercise figure and simulator figure use
different rotor numbering; the matrices in this section follow the exercise
figure, not the simulator figure.
Let $d>0$ be the centre-to-rotor distance and $a=d/\sqrt{2}$. The rotor
positions in the exercise figure are
\[
 \rho_1=(-a,a,0)^T,\quad \rho_2=(-a,-a,0)^T,\quad
 \rho_3=(a,-a,0)^T,\quad \rho_4=(a,a,0)^T.
\]
Define $u_i=\omega_i^2\geq0$, where $\omega_i$ is the speed magnitude in
each propeller's configured lifting direction, $c_f>0$ the lift coefficient,
and $c_d>0$ the reaction-torque coefficient. Then
\[
 f_i=c_fu_i,\qquad
 \rho_i\times (0,0,f_i)^T=(\rho_{iy}f_i,-\rho_{ix}f_i,0)^T.
\]
Counter-clockwise rotor motion produces clockwise body reaction torque,
so rotors 1 and 3 in (a) have negative yaw torque and rotors 2 and 4 have
positive yaw torque. With $w=(f,M_x,M_y,M_z)^T$, we obtain
\begin{equation}
 w=F_a u,\qquad
 F_a=\begin{bmatrix}
 c_f&c_f&c_f&c_f\\
 ac_f&-ac_f&-ac_f&ac_f\\
 ac_f&ac_f&-ac_f&-ac_f\\
 -c_d&c_d&-c_d&c_d
 \end{bmatrix}.
\end{equation}
The four scaled rows are mutually orthogonal and
\[
 \det F_a=16a^2c_f^3c_d=8d^2c_f^3c_d\ne0,
 \qquad \operatorname{rank}(F_a)=4.
\]
Thus all four wrench components can be locally adjusted independently
around a positive-speed hover, subject to actuator limits.

In configuration (b), swapping the propeller types and reversing the
directions of rotors 3 and 4 preserves upward lift but reverses their
yaw reaction torques. Therefore
\begin{equation}
 F_b=\begin{bmatrix}
 c_f&c_f&c_f&c_f\\
 ac_f&-ac_f&-ac_f&ac_f\\
 ac_f&ac_f&-ac_f&-ac_f\\
 -c_d&c_d&c_d&-c_d
 \end{bmatrix}.
\end{equation}
Its fourth row is $-c_d/(ac_f)$ times its second row. The first three
rows remain independent, hence
\[
 \boxed{\operatorname{rank}(F_b)=3,\qquad
 M_z=-\frac{c_d}{ac_f}M_x.}
\]
For example, $F_b(1,-1,1,-1)^T=0$. A pure yaw moment at zero roll moment
is impossible. At an instant of level hover with zero angular velocity,
keeping the position fixed requires zero roll acceleration and therefore
$M_x=0$, which also forces $M_z=0$. A prescribed yaw acceleration at that
instant cannot be realized. Consequently (b) cannot track an arbitrary
independently prescribed output trajectory $(x,y,z,\Psi)$; only trajectories
compatible with the moment constraint may be feasible. Changing coordinate
axes or wrench units cannot restore the missing rank.

\section*{B. Control with quadratic body drag}
Let $R\in SO(3)$ rotate body vectors into the world frame, $v$ be the
world velocity, and $v^b=R^Tv$. Write
\[
 D=\operatorname{diag}(0.1,0.1,0.2),\qquad
 h(v^b)=v^b\odot|v^b|,\qquad
 F_D^b=-D h(v^b),\qquad F_D^w=RF_D^b.
\]
Here $\odot$ means componentwise multiplication. This sign convention
opposes motion in every axis and satisfies
\[
 (v^b)^TF_D^b=-0.1|v_x^b|^3-0.1|v_y^b|^3-0.2|v_z^b|^3\leq0.
\]
With the lab's upward $z$ convention, the translational dynamics become
\[
 m\dot v=-mg e_3+fRe_3+F_D^w,\qquad e_3=(0,0,1)^T.
\]
Define $e_x=x-x_d$, $e_v=v-v_d$ and command the rotor force vector
\begin{equation}
 \boxed{A=-k_xe_x-k_ve_v+mg e_3+m\ddot x_d-\widehat F_D^w.}
\end{equation}
Using measured velocity and attitude, compute $R^Tv$, evaluate the
quadratic drag law in the body frame, then rotate the result by $R$.
Subtract this predicted world force when constructing $A$. For example,
with $R=I$ and $v=(2,0,0)^T$, drag is $(-0.4,0,0)^T$ N, so the
controller adds $+0.4$ N of forward force.

The compensated force changes both the desired thrust magnitude and
attitude. For $\|A\|>0$, choose
\[
 b_{3d}=\frac{A}{\|A\|},\quad b_{1c}=(\cos\Psi_d,\sin\Psi_d,0)^T,
 \quad b_{2d}=\frac{b_{3d}\times b_{1c}}{\|b_{3d}\times b_{1c}\|},
\]
\[
 R_d=[\,b_{2d}\times b_{3d}\quad b_{2d}\quad b_{3d}\,],\qquad
 \boxed{f=A^TRe_3.}
\]
With perfect attitude tracking, $fRe_3=A$. If the drag estimate is exact,
substitution into the dynamics gives
\[
 m\dot e_v=-k_xe_x-k_ve_v,
\]
so the modelled disturbance is cancelled rather than left for feedback
to counter through a persistent tracking error. With imperfect estimates,
the residual disturbance is $F_D^w-\widehat F_D^w$.

The geometric attitude controller retains its form:
\[
 e_R=\tfrac12(R_d^TR-R^TR_d)^\vee,\quad
 e_\Omega=\Omega-R^TR_d\Omega_d,
\]
\[
 M=-k_Re_R-k_\Omega e_\Omega+\Omega\times J\Omega
 -J(\widehat\Omega R^TR_d\Omega_d-R^TR_d\dot\Omega_d).
\]
If the full feedforward terms are used, $\Omega_d$ and $\dot\Omega_d$ must
come from the compensated $R_d$. The drag model supplied in the exercise
is a force at the centre of mass; it does not specify an extra body moment.

In the paper's downward $z$ convention, instead use
$A_p=-k_xe_x-k_ve_v-mg e_3+m\ddot x_d-\widehat F_D^w$,
$b_{3d}=-A_p/\|A_p\|$ and $f=-A_p^TRe_3$.
The implementation provides optional drag compensation, disabled by default
because the hypothetical exercise drag need not match the TESSE plant.

\section*{References}
\begin{enumerate}
\item MIT VNAV, \emph{Lab 3 Exercises (2023)},
\url{https://vnav.mit.edu/labs_2023/lab3/exercises.html}.
\item T. Lee, M. Leok, and N. H. McClamroch, \emph{Geometric Tracking
Control of a Quadrotor UAV on SE(3)}, CDC 2010, pp. 5420--5425,
\url{https://mleok.science/pdf/LeLeMc2010_quadrotor.pdf}.
\end{enumerate}
\end{document}
